\section{Command-line tools}

\subsection{Calibrating your own filters: \code{make\_filter\_profile.py}}

Amateur filters (Astrodon, Astronomik, Baader, Optolong RGB and
narrowband) are not in the SVO database, but a transmission curve is all
that is needed --- the zero points are computed, not measured:

\begin{verbatim}
python3 make_filter_profile.py Astrodon.G astrodon_g.dat nm
\end{verbatim}

reads a two-column wavelength/transmission file (0--1 or percent,
auto-detected; unit Angstrom, nm or um), and writes
\code{Astrodon.G\_Vega.xml} and \code{Astrodon.G\_AB.xml} into
\code{data/filters/} in full SVO format: complete VO annotation, all
characterising wavelengths, FWHM, effective width, the mean solar flux,
and a Vega zero point computed synthetically from the bundled CALSPEC Vega
spectrum --- the same convention SVO uses (validated to $\lesssim$1\%
against the shipped SVO profiles). Add the two printed lines to
\code{data/filters.list} and the filter appears in the selector on the
same photometric scale as the professional passbands.

\subsection{Importing template spectra: \code{add\_template.py}}

Template spectra can be two-column ASCII \emph{or} CALSPEC/STScI-style
FITS tables (\code{WAVELENGTH}/\code{FLUX}; Angstrom, nm or micron units
read from \code{TUNIT}). To add a single file or an entire directory ---
e.g.\ the whole CALSPEC archive folder:

\begin{verbatim}
python3 add_template.py current_calspec/
python3 add_template.py my_spectrum.fits "HD 12345" G2V
\end{verbatim}

For each spectrum the tool copies the file under \code{data/imported/},
computes the synthetic Vega $V$ magnitude (the catalogue \code{mv0}) and
$B-V$ against the shipped Bessell profiles, cleans CALSPEC filename
suffixes for the display name (\code{alpha\_lyr\_stis\_011} $\to$
\code{alpha\_lyr}), reports the $V$-band coverage, and appends the row to
\code{data/interpola.db.csv} --- the file the GUI template list is built
from. Press \textsf{Reload catalog} (or restart) and the new templates
appear in the selector. Surface-brightness and arbitrary-flux files
(planet atlas, supernova templates) are fully usable: the ETC rescales
every template to your entered magnitude, so only the spectral shape and
the computed \code{mv0} matter. Files that do not fully cover the $V$
band are flagged; use $V$ as the reference filter for those.

\subsection{Headless runs: \code{spetc\_batch.py}}

\begin{verbatim}
python3 spetc_batch.py examples/batch_photometry.json out/
\end{verbatim}

runs the same engines without the GUI from a JSON run description (two
commented examples ship in \code{examples/}) and writes the full result
CSV plus a self-contained \textbf{one-page HTML summary}: configuration,
sky, key numbers ($\sigma$(EW) for spectroscopy, saturation, maximum
unsaturated exposure), the stack plan, and an embedded S/N figure (S/N
versus exposure for photometry; S/N versus wavelength for spectroscopy).
Batch sky modes are \code{fixed\_ab} and \code{sqm}; the Moon and twilight
terms need the GUI's time machinery and are not applied, as the summary
footer states. Useful for scripting parameter studies and for archiving a
session plan.
